% 负责人：A / hycx233。数据准备与共享预处理的已验证事实。
\section{数据来源与准备}
\subsection{样本与标注}
使用课程提供的 Micro-C 接触矩阵和\texttt{标注数据.xlsx}。
原研究公开数据登记于 GEO 总系列 GSE272161\cite{geo}。
WT 使用 Micro-C 子系列 GSE272159 所列的两个 37\,$^\circ$C 附件\cite{geomicroc}，
突变数据实际输入取自课程提供的 \texttt{GSE272161\_RAW.tar}，
选择性解压 GSM8950761 至 GSM8950764 四个成员。
这四个样本也列于 GSE272159；本项目以课程配套文件为实际输入。
本阶段只读原材料，转换后的标注、基因表和样本表分别保存。

\begin{table}[htbp]
\centering
\caption{已核对的数据准备结果（任务 \#1）}
\begin{tabularx}{\linewidth}{l r X}
\toprule
产物 & 记录数 & 内容 \\
\midrule
结构表 & 344 & OPCID 68 条、CHIN 250 条、CHID 26 条 \\
基因表 & 4,516 & 4,451 个不同名称；保留同名但位置不同的记录 \\
样本表 & 6 & WT、$\Delta stpA$、$\Delta hns\Delta stpA$ 各两个重复 \\
\bottomrule
\end{tabularx}
\end{table}

六个实际矩阵的染色体名均为 \texttt{NC\_000913.3}，长度为 4,641,652 bp，
分辨率为 10 bp，每个样本有 464,166 个 bin；计数为 \texttt{int32}，
采用上三角对称存储，均无预存的 \texttt{weight} 列。
因此共享预处理显式读取原始 counts，不假定平衡权重已经存在。

\subsection{坐标口径与追溯}
下游统一使用染色体名 \texttt{NC\_000913.3} 和 0-based 左闭右开区间。
Excel 中的 \texttt{MG1655} 作为同一染色体的别名。

Table 1 的基因坐标与 NCBI 的 1-based 闭区间一致，转换时仅起点减一，终点和链方向不变。
例如 \textit{thrL} 的源范围为 $190\ldots255$，转换后为 $[189,255)$；
\textit{hns} 的源范围为 $1292509\ldots1292922$，转换后为 $[1292508,1292922)$，
仍为负链\cite{thrl,hns}。

结构表沿用作者接触图坐标的数值，不减一。
作者分析代码直接以坐标整除 25 进行分箱\cite{authorcode}。
将课程标注按该规则转为作者的 BED/BEDPE 格式，
68 条 OPCID、250 条 CHIN（区间及中心）、26 条 CHID 均匹配。
原文未明确说明 Excel 结构端点逐碱基的开闭含义，
因此将其解释为 $[start,end)$ 是本项目的统一接口约定。
CHIN 的 Center 复制原始字段，所有记录保留来源工作表、行号与原坐标。

\subsection{数据准备检查}
结构表与基因表已逐行回查 Excel，检查计数、唯一标识、坐标范围、长度和 CHIN 原 Center。
六份矩阵已核对元数据，解压文件可复用，原始矩阵和压缩包不改写。
还在 WT 两个重复中各提取了三类结构的真实窗口，作为初始 CPU 读取与定位预览。
这批初始预览显示统一色标下的 $\log(1+\mathrm{raw\ counts})$，
仅用于核对读取位置，不构成条件差异证据；正式共享输入使用下述归一化。

运行入口为 \path{scripts/prepare_annotations.py}、
\path{scripts/prepare_samples.py} 和 \path{scripts/preview_structures.py}。
详细命令、表字段和代表预览保存在仓库的\texttt{docs/数据准备说明.md}。

\subsection{共享归一化与有效掩码}
采用流式聚合，将原始 10 bp 上三角唯一接触计数相加至 100 bp 工作矩阵，
读取局部窗口时补成对称矩阵，不展开全基因组稠密数组。
令 $C_{ij}$ 为工作矩阵计数，样本总接触量和深度归一化强度分别为
\[
  T=\sum_{i\leq j}C_{ij},\qquad I_{ij}=\frac{10^6 C_{ij}}{T}.
\]
$T$ 对每个接触计数一次；轨道的边缘信号取对称矩阵行和，
非对角计数加到两个端点，对角计数只加一次。
不同条件的强度比较使用 $I$，不把每个窗口独立缩放到同一均值。

有效 bin 要求边缘信号大于零且宽度完整为 100 bp。
对 $n$ 个工作 bin，定义环状距离 $d(i,j)=\min(|i-j|,n-|i-j|)$。
令 $P_d$ 为距离 $d$ 上两个端点都有效的无序 bin 对集合，包含零接触对，
则形态背景、O/E 和统一变换为
\[
  E_d=\frac{\sum_{(i,j)\in P_d}C_{ij}}{|P_d|},\qquad
  R_{ij}=\frac{C_{ij}}{E_{d(i,j)}},\qquad
  F_{ij}=\frac{\log(1+\min(R_{ij},10))}{\log(11)}.
\]
O/E 仅在 $E_d>0$ 且端点有效时计算，形态另掩去主对角线。
强度保留对角并有独立掩码。无效值以零存储，不能据此解释为接触减弱；
在同一位置比较多个样本时，要使用这些样本有效掩码的交集。
全局深度和距离背景不使用结构标签或模型指标，截断值在模型评价前固定。

\subsection{环状窗口与模型输入}
窗口以结构中心为中心，跨度 24 kb。跨起点时按真实基因组长度 $L=4641652$
拆成首尾两段，读取四个矩阵块，保留两段之间的非对角接触。
100 bp bin 与窗口端点不一定对齐，局部矩阵通常有 240--242 格。
生成 $128\times128$ 输入时，以真实 bin 区间与目标格的重叠面积作权重，
只平均有效面积，同时保存模型掩码与有效面积比例。

末尾 52 bp bin 的原始计数保留，但不参与强度及形态统计。
距离背景使用 100 bp 网格上的环状距离近似，跨起点距离最多相差 48 bp；
真实窗口坐标和面积权重仍使用实际 $L$，不会把染色体补长。

\subsection{分组与导出}
对已知结构的完整 24 kb 窗口求环状空间重叠连通分量，
并检查实际读取的 100 bp 边缘 bin。同一分量的结构、重复及后续增强使用同一集合。
固定种子为 20260918，按类别数与总量分配约 70\%/15\%/15\%，不使用模型结果选择划分。
344 条结构形成 59 个分量，最大分量 20 条；没有排除任何已知结构。

\begin{table}[htbp]
\centering
\caption{共享划分：结构位置与 WT 两个重复的输入窗口}
\begin{tabular}{l r r r r r}
\toprule
集合 & OPCID & CHIN & CHID & 位置总数 & 两重复窗口 \\
\midrule
train & 48 & 175 & 18 & 241 & 482 \\
val & 10 & 37 & 4 & 51 & 102 \\
test & 10 & 38 & 4 & 52 & 104 \\
\bottomrule
\end{tabular}
\end{table}

分类导出 688 个输入。另导出 24 个随机背景位置、共 48 个重复窗口，用于发现模块接通入口；
背景位置不与已知结构注释相交，窗口不得连接不同集合，
与已有窗口重叠时继承其集合。六样本检查另有 24 个示例窗口，其中跨起点检查标记为
\texttt{probe}，不参与训练、验证或测试。
背景和跨起点 \texttt{probe} 的标签为 $-1$，不能混入三分类训练。

\subsection{实际运行与验证}
2026 年 9 月 20 日已在 CPU 上用六份真实数据跑通任务 \#2。
首次缓存生成约耗时 253 秒，复用缓存导出窗口及绘图约 23 秒；
这些是本次运行记录，不是其他机器的耗时保证。
九项针对性检查通过，六个样本聚合前后的总计数一致。
WT rep1 的五个普通及首尾窗口另与原始 10 bp 局部像素的独立聚合结果逐元素一致。
跨集合共享的实际 100 bp bin 数为零。
688 个分类输入已由分类 PR \#13 的实际读取代码加载，数组、标签与 CSV 行序一致，
且保留提供的划分；背景与示例输入也通过有限数、对称性、标签、行序和坐标检查。
\textbf{本阶段没有运行模型训练，不报告模型性能。}

安装仓库依赖并准备原始样本后，从仓库根目录执行：
\begin{verbatim}
OPENBLAS_NUM_THREADS=1 python scripts/prepare_data.py
\end{verbatim}
详细方法、验证记录和运行时的代码指纹保存在以下位置：
\begin{flushleft}
\small
方法说明：\texttt{docs/共享预处理说明.md}\\
验证记录：\path{docs/results/preprocessing/validation.json}\\
分组摘要：\path{docs/results/preprocessing/splits.json}\\
运行记录：\path{outputs/preprocessing/default/run_*.json}
\end{flushleft}
